% Part: history
% Chapter: biographies 
% Section: alan-turing 

\documentclass[../../../include/open-logic-section]{subfiles}

\begin{document}

\olfileid{his}{bio}{tur} 

\olsection{ایلن ٹیورنگ}

ایلن ٹیورنگ 23 جون 1912 نوں مَیڈا ویل، لندن وچ جمیا۔ اوہنوں نظری کمپیوٹر سائنس دا بانی منّیا جاندا اے۔ طبیعی سائنس تے ریاضی وچ ٹیورنگ دی دلچسپی نِکّی عمر توں شروع ہوئی۔ پر اوہ جیہڑے سکولاں وچ پڑھیا، اوتھے اوہدی دلچسپی دے مضمون گھٹ پڑھائے جاندے سن۔ اوتھے ادب تے کلاسیکی لکھتاں اُتے زور سی۔ ایس لئی سکول وچ اوہدی کارکردگی ماڑی رہی تے کئی استاداں نے اوہنوں جھاڑ پائی۔

\olphoto{turing-alan}{ایلن ٹیورنگ}

ٹیورنگ نے کیمبرج دے کنگز کالج وچ پہلی ڈگری لئی ریاضی پڑھی۔ 1936 وچ ٹیورنگ نے اوہ شے بنائی جیہنوں ہُن ٹیورنگ مشین آکھدے نیں۔ اوہدی کوشش سی کہ حساب جوگ فنکشن دی عین تعریف دے سکے تے ثابت کرے کہ فیصلہ مسئلے دا فیصلہ نہیں ہو سکدا۔ الونزو چرچ ایس نتیجے تک پہلاں پہنچ گیا۔ اوہنے اپنے لیمبڈا کیلکولس نال ایہ نتیجہ ثابت کیتا سی۔ ٹیورنگ دا مقالہ فیر وی چرچ دے نتیجے دا حوالہ دے کے چھپیا۔ چرچ نے ٹیورنگ نوں پرنسٹن سَدیا۔ اوہنے 1936 توں 1938 تک اوتھے وقت گزاریا تے چرچ دی نگرانی تھلّے ڈاکٹریٹ لئی۔

منطق وچ دلچسپی دے باوجود ٹیورنگ دی طبیعی سائنس وچ پہلی دلچسپی قائم رہی۔ دوجی عالمی جنگ وچ بلیچلی پارک دے برطانوی خفیہ رمزاں دی پڑتال والے محکمے وچ کم کردیاں اوہدی عملی مہارت کم آئی۔ جرمن بحری فوج دے پیغاماں لئی ورتے جاندے اینگما رمز نوں توڑن وچ ٹیورنگ دا مرکزی کردار سی۔ شماریات تے خفیہ رمزاں وچ ٹیورنگ دی مہارت تے الیکٹرانک مشیناں دی آمد نال ٹیم «بومب» ناں دی رمز کھولھن والی مشین بنا کے رمز توڑ سکی۔ اوہدے خیالات نے دنیا دا پہلا پروگرام کیتا جا سکّن والا الیکٹرانک کمپیوٹر، کولاسس، بناؤن وچ وی مدد دتی۔ اوہ وی بلیچلی پارک وچ جرمن لورینز رمز توڑن لئی ورتیا گیا۔

ٹیورنگ ہم جنس پسند سی۔ فیر وی 1942 وچ اوہنے بلیچلی پارک دی اپنی ساتھی جون کلارک نوں ویاہ دی پیشکش کیتی۔ پر بعد وچ اوہنے منگنی توڑ دتی تے کلارک نوں دسیا کہ اوہ ہم جنس پسند اے۔ اوہدی زندگی وچ کئی محبوب سن، بھاویں اوس ویلے برطانیہ وچ ہم جنس جنسی عمل جرم سن۔ 1952 وچ ٹیورنگ دے اوس ویلے دے محبوب دے اک دوست نے ٹیورنگ دے گھر چوری کیتی۔ پولیس رپورٹ لکھاؤندیاں ٹیورنگ نے ہم جنس تعلق دا اقرار کیتا، کیوں جے اوہنوں لگدا سی کہ سرکار ہم جنس جنسی عملاں نوں قانونی کرن والی اے۔ ایہ گل سچ نہیں سی، تے اوہدے اُتے سنگین بے حیائی دا الزام لگایا گیا۔ جیل جان دی بجائے ٹیورنگ نے ہارمون والا علاج چُنیا جیہڑا جنسی خواہش گھٹاؤندا سی۔ ٹیورنگ 8 جون 1954 نوں مردہ لبھیا۔ موت سائینائیڈ دی ودھ مقدار نال ہوئی، غالباً خودکشی سی۔ 2013 وچ ملکہ الزبتھ~دوجی نے اوہنوں شاہی معافی دتی۔

\begin{reading}
ایلن ٹیورنگ دی تفصیلی سوانح عمری لئی \citet{Hodges2014} ویکھو۔ ٹیورنگ دی زندگی تے کم توں \emph{Breaking the Code} ناں دا ڈراما بنیا۔ 1996 وچ ایہ ٹی وی لئی بنیا، جیہدے وچ ڈیرک جیکوبی نے ٹیورنگ دا کردار کیتا۔ اکیڈمی انعام لئی نامزد فلم \emph{The Imitation Game} وچ بینیڈکٹ کمبربیچ تے کیرا نائٹلی نیں۔ ایہ وی ٹیورنگ دی زندگی تے بلیچلی پارک وچ اوہدے وقت اُتے ڈھیلے ڈھالے طور اُتے مبنی اے \citep{Imitation2014}۔

\citet{Radiolab2012} وچ ٹیورنگ دی زندگی تے کم بارے کئی پوڈکاسٹ نیں۔ بی بی سی ہورائزن دی دستاویزی فلم \emph{The Strange Life and Death of
  Dr.~Turing} انٹرنیٹ اُتے ویکھی جا سکدی اے \citep{Sykes1992}۔ \citep{Theelen2012} وچ لیگو دی کم کردی ٹیورنگ مشین دی مختصر ویڈیو اے۔ ایہ 2012 وچ ٹیورنگ دے جنم دے سو سال پورے ہون اُتے اوہدی یاد وچ بنائی گئی۔

ٹیورنگ مشیناں تے فیصلہ مسئلے اُتے ٹیورنگ دا اصل مقالہ \citet{Turing1937} اے۔
\end{reading}

\end{document}
